%% kapa: exam=1 round=20 year=2009 period=1 form=A subject=economics subject_ko=경제원론 seq=21 number=61
%% source: Q-Net 감정평가사 기출문제 (cst00309) artlSeq=5024178 file=extracted/1008/5024178/2009년 제20회 감정평가사 1차시험 기출문제/1교시.hwp
%% conversion: kordoc 4.24.51 -> tools/gen1.py (auto); review=proofread-sonnet; flags=none
\begin{question}{61}
소비자 A는 두 재화 $X$ 와 $Y$ 만을 소비하고 있다. 그의 무차별곡선은 $y= \frac{ U } { x+6 }$로 표현된다. 다음 설명 중 옳은 것은? (단, $U$ 는 효용수준, $x$ 는 $X$ 재의 소비량, $y$ 는 $Y$ 재의 소비량이다.)
\begin{choices}
  \item $X$ 재 1단위와 $Y$ 재 1단위를 소비하면 A의 효용수준은 5이다.
  \item $X$ 재 2단위와 $Y$ 재 2단위를 소비하면 A의 효용수준은 10이다.
  \item $X$ 재 15단위와 $Y$ 재 10단위를 소비하는 것보다 $X$ 재 10단위와 $Y$재 15단위를 소비하는 것을 선호한다.
  \item $X$ 재 8단위와 $Y$ 재 9단위를 소비하는 것보다 $X$ 재 9단위와 $Y$ 재 8단위를 소비하는 것을 선호한다.
  \item $X$ 재 8단위와 $Y$ 재 9단위를 소비하는 것과 $X$ 재 9단위와 $Y$ 재 8단위를 소비하는 것은 무차별하다.
\end{choices}
\end{question}
\answer{3}
